\phantomsection
\addcontentsline{toc}{section}{About Part IB}
\section*{About Part IB}

If Part IA is about technique, Part IB is about abstraction. This is the year in which
the objects you have been computing with are replaced by the axioms they satisfy, and
it is by some distance the largest single jump in the degree. It is also the year with
the most courses in it -- fifteen chapters here, out of a longer list on offer --
because the second year is where a mathematician's tastes are supposed to start
declaring themselves.

The clearest example of the shift is Linear Algebra. Vectors and Matrices spent a year
in $\R^3$ and $\C^n$; Linear Algebra deletes the coordinates and starts from a vector
space over an arbitrary field, and suddenly the dual space, the trace, the determinant
and the Jordan normal form all become statements about linear maps rather than
statements about arrays of numbers. The same thing happens to analysis. Analysis and
Topology takes the $\epsilon$--$N$ arguments of Analysis I and asks what they actually
needed; the answer is a metric, and then not even that, and by the end of the chapter
you are proving compactness results in a topological space with no numbers in sight.
Groups, Rings and Modules performs the same trick on Groups, generalising the integers
and the polynomials into a single theory of factorisation, and then generalising vector
spaces into modules so thoroughly that the classification of finitely generated modules
over a Euclidean domain hands you back both the structure of finite abelian groups and
the Jordan normal form as corollaries.

That triple -- Linear Algebra, Analysis and Topology, Groups, Rings and Modules -- is
the spine of the year, and almost every Part II course leans on at least one of them.

Alongside the abstraction there is theory-building of a different kind. Complex Analysis
is the outstanding course of the second year, and possibly of the degree: a short list
of hypotheses about a function of a complex variable turns out to force an
extraordinary amount of structure, and Cauchy's theorem propagates through the chapter
until residues, the argument principle and Rouch\'e's theorem all fall out of it.
Complex Methods covers overlapping ground from the other direction, caring about
contour integrals you can actually evaluate and conformal maps you can actually draw;
reading the two chapters side by side is instructive precisely because they disagree
about what a proof is for. Geometry, meanwhile, does for surfaces what Analysis and
Topology does for spaces, working up through the two fundamental forms to the
Theorema Egregium and the Gauss--Bonnet theorem, which is the first result in the book
where a local computation and a global topological invariant are forced to agree.

The applied courses inherit Part IA's machinery and finally put it to work. Vector
Calculus's integral theorems become Maxwell's equations in Electromagnetism and the
transport theorems in Fluid Dynamics. Methods is the workhorse of the year, and worth
singling out: Fourier series, Sturm--Liouville theory, distributions and Green's
functions form a toolkit that reappears in nearly every applied course afterwards.
Variational Principles reframes mechanics in terms of stationary actions, which is the
language Principles of Quantum Mechanics and Statistical Physics will later use, and
Quantum Mechanics itself is in large part an exercise in solving the equations that
Methods taught you to solve.

Four courses here sit slightly apart. Statistics and Markov Chains take Part IA
Probability seriously enough to build real theory on it -- estimators, confidence
intervals and hypothesis tests in one case, recurrence, invariant distributions and
convergence to equilibrium in the other. Optimisation, with its convexity, duality and
simplex method, and Numerical Analysis, with its interpolation, quadrature and
stability theory, are the two most computational chapters in the book. None of them is
a prerequisite for a great deal, but Optimisation and Markov Chains both feed into
Stochastic Financial Models, and Numerical Analysis shares more than you would expect
with Topics in Analysis.

The character of the year, in short, is that objects stop being defined by what they
are and start being defined by what they do. It is uncomfortable for a while, and then
it becomes the only way you want to think.
